\chapter{Optimizers and the Training Loop from First Principles}\label{ch:optim}

Part~II built a model. It cannot do anything until something moves its weights, and that
something is a surprisingly large system: a sampler that pulls windows out of a token file
holding billions of entries, a numerically careful loss, an optimizer that keeps two extra
copies of every parameter, a learning-rate schedule that must know in advance when the run ends,
a clipping step that guards against rare bad batches, a checkpointer that has to survive a
crash, and a logger that records enough for you to work out later what happened. Each
piece is a few dozen lines. Getting all of them right, and fast, separates a training
script from a training \emph{system}.

This chapter builds every one of those pieces from first principles: SGD, momentum, AdaGrad,
RMSProp, Adam and AdamW derived and implemented as \code{torch.optim.Optimizer} subclasses;
Lion and Muon as modern alternatives; warmup, cosine and warmup--stable--decay schedules;
global-norm clipping; a resumable memory-mapped batch sampler; atomic, bit-exact
checkpoints; and a small experiment-management layer. All of it lives in
\ident{llmsys/optim.py} and \ident{llmsys/explog.py}. \Cref{ch:pretrain} then puts these pieces
into the GPT-2 pretraining loop and adds mixed precision, activation checkpointing and MFU.
Throughout, we treat the optimizer step as what it is on a GPU: an elementwise,
memory-bound kernel whose cost we can predict to within a few percent.

\begin{figure}[t]
\centering
\begin{tikzpicture}[
  st/.style={draw, rounded corners=3pt, minimum height=1.0cm, text width=2.15cm, align=center,
             font=\small, line width=0.6pt},
  num/.style={circle, fill=sysblue, text=white, font=\scriptsize\bfseries, inner sep=1.5pt},
  arr/.style={-{Stealth[length=2mm]}, line width=0.6pt}]
\node[st, fill=fdata!18, draw=keygreen] (a) at (0,0) {Loss \& the\\Optimizer API};
\node[st, fill=fcompute!18, draw=sysblue] (b) at (3.25,0) {SGD $\to$ Adam\\$\to$ AdamW};
\node[st, fill=fmemory!20, draw=napkin] (c) at (6.5,0) {Cost of a step:\\bytes \& FLOPs};
\node[st, fill=fcompute!18, draw=sysblue] (d) at (9.75,0) {Schedules,\\clipping, Muon};
\node[st, fill=fdata!18, draw=keygreen] (e) at (0,-2.1) {Batches from\\memmapped tokens};
\node[st, fill=fmemory!20, draw=napkin] (f) at (3.25,-2.1) {Atomic,\\bit-exact\\checkpoints};
\node[st, fill=fcomm!16, draw=labpurple] (g) at (6.5,-2.1) {Configs, logs,\\sweeps};
\node[st, fill=warnred!12, draw=warnred] (h) at (9.75,-2.1) {Debugging\\the loop};
\foreach \n/\k in {a/1,b/2,c/3,d/4,e/5,f/6,g/7,h/8} \node[num] at (\n.north west) {\k};
\draw[arr] (a)--(b); \draw[arr] (b)--(c); \draw[arr] (c)--(d);
\draw[arr] (d.south) -- ++(0,-0.35) -| (e.north);
\draw[arr] (e)--(f); \draw[arr] (f)--(g); \draw[arr] (g)--(h);
\end{tikzpicture}
\caption{Roadmap of this chapter. Stages 1--4 are the mathematics and cost of the update rule;
stages 5--8 are the plumbing that turns an update rule into a training run you can trust.}
\label{fig:optim:roadmap}
\end{figure}

% =====================================================================
\section{The training loop as a system}\label{sec:optim:loop}\covers{CS336 L2}

Every training step does the same seven things (\cref{fig:optim:loop}): fetch a batch,
compute logits, reduce them to a scalar loss, backpropagate, possibly clip the gradient, update
the parameters (and the optimizer's state), and occasionally evaluate, log and checkpoint. In
code:

\begin{lstlisting}[caption={The skeleton every training loop in this book follows.},label={lst:optim:skeleton}]
for step in range(start_step, total_steps):
    set_lr(opt, schedule(step))               # 1. schedule
    x, y = next(batches)                      # 2. data (B, T) int64 on device
    logits = model(x)                         # 3. forward  (B, T, V)
    loss = cross_entropy(logits, y)           # 4. loss     scalar
    opt.zero_grad(set_to_none=True)
    loss.backward()                           # 5. backward
    gnorm = clip_grad_norm_(model.parameters(), 1.0)   # 6. clip
    opt.step()                                # 7. update p and its state
    if step % log_every == 0:
        logger.log(step, loss=loss.item(), grad_norm=gnorm)
    if (step + 1) % ckpt_every == 0:
        ckpt.save(step + 1, model, opt, batches)
\end{lstlisting}

Two properties of this loop matter more than any single line in it. First, it must be
\emph{resumable}: a job on a shared cluster will be preempted, so the state that defines
``where we are'' (weights, optimizer moments, step, data position, random-number generators)
must be saved and restored completely. Second, it must be \emph{observable}: when a run goes
wrong after twelve hours you need the gradient norm, learning rate and loss at every step.
\Cref{sec:optim:ckpt,sec:optim:explog} build both properties.

\begin{figure}[t]
\centering
\resizebox{\linewidth}{!}{%
\begin{tikzpicture}[
  box/.style={draw, rounded corners=2pt, minimum height=0.9cm, align=center, font=\small,
              line width=0.6pt, text width=2.1cm},
  arr/.style={-{Stealth[length=2mm]}, line width=0.6pt},
  side/.style={draw, dashed, rounded corners=2pt, align=center, font=\footnotesize, text width=2.2cm}]
\node[box, fill=fdata!18, draw=keygreen] (data) at (0,0) {token memmap\\on disk};
\node[box, fill=fdata!18, draw=keygreen] (samp) at (2.9,0) {random-window\\sampler};
\node[box, fill=fcompute!18, draw=sysblue] (fwd) at (5.6,0) {forward\\$(B,T,V)$ logits};
\node[box, fill=fcompute!18, draw=sysblue] (loss) at (8.4,0) {fp32\\cross-entropy};
\node[box, fill=fcompute!18, draw=sysblue] (bwd) at (11.2,0) {backward\\grads};
\node[box, fill=fmemory!20, draw=napkin] (clip) at (11.2,-1.9) {global-norm\\clip};
\node[box, fill=fmemory!20, draw=napkin] (opt) at (8.4,-1.9) {optimizer step\\$p, m, v$};
\node[box, fill=fcomm!16, draw=labpurple] (sch) at (5.6,-1.9) {LR schedule\\$\eta(t)$};
\node[side] (ck) at (2.8,-1.9) {checkpoint:\\weights, $m,v$, step,\\sampler, RNG};
\node[side] (lg) at (0,-1.9) {JSONL log:\\step, wall, loss,\\lr, $\|g\|$};
\draw[arr] (data)--(samp);
\draw[arr] (samp)--(fwd);
\draw[arr] (fwd)--(loss); \draw[arr] (loss)--(bwd); \draw[arr] (bwd)--(clip);
\draw[arr] (clip)--(opt); \draw[arr] (sch)--(opt);
\draw[arr, dashed] (opt)--(ck); \draw[arr, dashed] (samp)--(ck);
\draw[arr, dashed] (ck)--(lg);
\draw[arr] (opt.north) |- ([yshift=-0.5cm]fwd.south) node[pos=0.75,above,font=\scriptsize]{new $p$} -- (fwd.south);
\end{tikzpicture}}
\caption{Dataflow of one training step. Solid arrows run every step; dashed arrows run every
few hundred steps. The checkpoint must capture everything on the left of the dashed line,
including the sampler's random state, or a resumed run will not reproduce the original.}
\label{fig:optim:loop}
\end{figure}

\begin{systemslens}[Where the time goes in the loop]
On a GPU the forward and backward passes are matrix multiplies and dominate the step: roughly
$6N$ FLOPs per token for an $N$-parameter model (\cref{ch:gpt}). Everything else in
\cref{lst:optim:skeleton} is elementwise or host-side work and should cost a few percent:
\begin{itemize}
  \item the optimizer step reads and writes every parameter, gradient and moment once
        (memory-bound, \cref{sec:optim:cost});
  \item clipping needs one reduction over all gradients and a host synchronisation if you turn
        the norm into a Python float;
  \item \code{loss.item()} and \code{float(gnorm)} each force the CPU to wait for the GPU, so call
        them only on logging steps or you serialise the asynchronous pipeline of \cref{ch:pytorch};
  \item batch sampling runs on the host and must overlap with GPU work (\cref{sec:optim:data}).
\end{itemize}
If a profile shows any of these at more than 5--10\% of step time, it is a bug, not a
fundamental cost.
\end{systemslens}

% =====================================================================
\section{The objective, numerically}\label{sec:optim:loss}\covers{CS336 A1 \S4.1}

The language-modelling loss is the average negative log-likelihood of the next token,
\begin{equation}
\ell(\theta) = -\frac{1}{BT}\sum_{b,t}\log \softmax(\vect{o}_{b,t})_{y_{b,t}}
= \frac{1}{BT}\sum_{b,t}\Bigl(\operatorname{logsumexp}(\vect{o}_{b,t}) - o_{b,t,y_{b,t}}\Bigr),
\label{eq:optim:ce}
\end{equation}
where $\vect{o}_{b,t}\in\R^{V}$ are logits. \Cref{sec:pretrain:loss} derives this and gives a
from-scratch implementation (\cref{lst:pretrain:ce}). Three numerical rules make it safe:

\begin{enumerate}
  \item \textbf{Never form $\exp(o)$ of raw logits.} With bf16 logits of magnitude 90,
        $e^{90}\approx 10^{39}$ overflows fp32. Subtract the row maximum first:
        $\operatorname{logsumexp}(\vect{o}) = m + \log\sum_u e^{o_u - m}$ with $m=\max_u o_u$.
        The same trick makes softmax stable, and it is the identity behind FlashAttention's
        online softmax (\cref{sec:attention:flash}).
  \item \textbf{Cancel $\log$ against $\exp$ symbolically.} Computing $\softmax$ then $\log$
        loses precision for tiny probabilities (and gives $\log 0=-\infty$); the right-hand form of
        \cref{eq:optim:ce} never materialises a probability.
  \item \textbf{Do the reduction in fp32.} The $(B,T,V)$ logits tensor is the largest activation
        in a small model, and averaging $BT$ terms in bf16 loses about three significant digits.
        Upcast the logits (or let a fused kernel do it) before the logsumexp.
\end{enumerate}

\begin{keyidea}[Know your loss at initialisation]
A freshly initialised model with small output weights predicts an almost uniform distribution
over $V$ tokens, so the first loss should be close to $\ln V$: about 6.24 for $V=512$ and
10.83 for $V=50{,}257$. A first loss of 30 means the logits are too large (initialisation or a
missing final norm); a first loss of exactly $\ln V$ to six digits often means the output is
constant (dead network, zeroed head). Perplexity $e^{\ell}$ starts near $V$ and should fall
quickly.
\end{keyidea}

The vocabulary dimension also makes the loss a \emph{memory} problem: the logits for
$B\cdot T=16{,}384$ tokens and $V=151{,}936$ (Qwen3) are 2.5\,G elements, or 10\,GB in fp32.
Production code therefore computes the loss in chunks over the sequence, or uses fused
``linear + cross-entropy'' kernels that never store the full logits. \Cref{ch:distill} implements
chunked cross-entropy for long reasoning traces.

% =====================================================================
\section{Writing an optimizer in PyTorch}\label{sec:optim:api}\covers{CS336 A1 \S4.2}

PyTorch's \code{torch.optim.Optimizer} base class does three jobs for us. It stores
\emph{parameter groups} (lists of tensors that share hyperparameters), it keeps a per-tensor
\emph{state} dictionary (\code{self.state[p]}), and it serialises both through
\code{state\_dict()}/\code{load\_state\_dict()}. A subclass only has to implement
\code{\_\_init\_\_} (pass the parameters and a \code{defaults} dictionary to \code{super()}) and
\code{step}. \Cref{lst:optim:sgd} is SGD with momentum, written the way PyTorch defines it so
that our test can compare against \code{torch.optim.SGD} to six digits.

\begin{lstlisting}[caption={SGD with heavy-ball or Nesterov momentum (from \texttt{llmsys/optim.py}).},label={lst:optim:sgd}]
class SGD(Optimizer):
    def __init__(self, params, lr=1e-2, momentum=0.0, nesterov=False, weight_decay=0.0):
        super().__init__(params, dict(lr=lr, momentum=momentum, nesterov=nesterov,
                                      weight_decay=weight_decay))

    @torch.no_grad()                        # the update itself is not part of the graph
    def step(self, closure=None):
        loss = None
        if closure is not None:             # API contract: optional loss re-evaluation
            with torch.enable_grad():
                loss = closure()
        for group in self.param_groups:     # each group may have its own lr / wd
            lr, mu, wd = group["lr"], group["momentum"], group["weight_decay"]
            for p in group["params"]:
                if p.grad is None:          # frozen or unused parameter
                    continue
                g = p.grad if wd == 0 else p.grad.add(p, alpha=wd)
                if mu != 0:
                    st = self.state[p]      # persistent, saved by state_dict()
                    if "momentum_buffer" not in st:
                        st["momentum_buffer"] = g.clone()
                    else:
                        st["momentum_buffer"].mul_(mu).add_(g)
                    buf = st["momentum_buffer"]
                    g = g.add(buf, alpha=mu) if group["nesterov"] else buf
                p.add_(g, alpha=-lr)        # in-place: no new tensor, no graph
        return loss
\end{lstlisting}

Several details in this listing are easy to get wrong:
\begin{itemize}
  \item \textbf{In-place updates under \code{no\_grad}.} \code{p.add\_} mutates the parameter's
        storage, so every reference to it (including one inside a tied module) sees the change.
        Writing \code{p = p - lr*g} would only rebind a local name. Some older code writes to
        \code{p.data}, which also works but bypasses autograd's version counter; the decorator is
        the modern idiom.
  \item \textbf{State lives in \code{self.state[p]}, keyed by the tensor.} That is what makes
        \code{state\_dict()} work: PyTorch maps tensors to integer indices in group order, so a
        checkpoint can be loaded into a freshly built model \emph{if parameters are created in the
        same order}. Reordering modules between save and load scrambles optimizer state silently.
  \item \textbf{Parameter groups are how you vary hyperparameters.} The usual LLM recipe puts
        matrices in a group with weight decay and gains/biases in a group without it
        (\code{decay\_param\_groups} in the library), and schedules write \code{group["lr"]} on
        every step (\cref{sec:optim:sched}).
  \item \textbf{\code{zero\_grad(set\_to\_none=True)}} frees the gradient tensors rather than
        filling them with zeros. That saves a memset kernel per parameter, and peak memory between steps.
\end{itemize}

\begin{pitfall}[A learning rate that is ``too small'' can look like a bug]
On a toy quadratic you can watch three regimes as you scale the learning rate up by factors of 10:
the loss falls slowly, falls fast, oscillates with growing amplitude, and finally diverges to
\code{inf} within a few steps. For gradient descent on $\tfrac12\lambda x^2$ the boundary is
exactly $\eta\lambda = 2$: the update multiplies $x$ by $(1-\eta\lambda)$, whose magnitude
exceeds 1 beyond that point. A real network has a whole spectrum of curvatures $\lambda_i$, and
the largest one (the ``sharpness'') sets the stability limit; this is the origin of the folk
wisdom that the best learning rate sits just below the ``edge of stability''
(\cref{sec:optim:debug}).
\end{pitfall}

% =====================================================================
\section{From SGD to Adam and AdamW}\label{sec:optim:adam}\covers{CS336 L2; CS336 A1 \S4.3}

Each classic optimizer adds one idea to SGD. Writing $g_t$ for the gradient at step $t$:

\begin{description}[leftmargin=1.2em, style=nextline]
  \item[Momentum.] $b_t = \mu b_{t-1} + g_t$, $\theta_t = \theta_{t-1} - \eta b_t$. Averaging
    gradients over roughly $1/(1-\mu)$ steps damps oscillation across steep directions and
    accumulates speed along shallow ones. Nesterov's variant evaluates the step ``ahead'' by using
    $g_t + \mu b_t$.
  \item[AdaGrad.] Keep $s_t = s_{t-1} + g_t^2$ (elementwise) and step
    $-\eta\, g_t/\sqrt{s_t+\epsilon}$. Each coordinate gets its own learning rate, inversely
    proportional to its accumulated gradient size. Because $s_t$ only grows, the effective step size
    decays to zero, which is good for convex problems and bad for long non-convex training.
  \item[RMSProp.] Replace the sum with an exponential moving average,
    $v_t = \beta_2 v_{t-1} + (1-\beta_2) g_t^2$, so the per-coordinate scale tracks the
    \emph{recent} gradient magnitude.
  \item[Adam] (Kingma and Ba, 2015) is RMSProp with momentum on the numerator,
    $m_t = \beta_1 m_{t-1} + (1-\beta_1) g_t$, plus a bias correction.
\end{description}

\paragraph{Why Adam needs bias correction.} With $m_0=0$, unrolling the recursion gives
$m_t = (1-\beta_1)\sum_{k=1}^{t}\beta_1^{t-k} g_k$. If the gradients have a stationary mean
$\bar g$, then $\E[m_t] = (1-\beta_1)\bar g \sum_{k=1}^{t}\beta_1^{t-k} = (1-\beta_1^t)\,\bar g$.
The estimate is biased towards zero early in training, by a factor of 10 at $t=1$ when
$\beta_1=0.9$. Dividing by $(1-\beta_1^t)$ removes the bias, and likewise $(1-\beta_2^t)$ for
$v_t$, where the effect is larger: with $\beta_2=0.999$ the uncorrected second moment would be
1000 times too small at step 1, and the first steps would be enormous. The full update is
\begin{equation}
\hat m_t = \frac{m_t}{1-\beta_1^t},\qquad \hat v_t = \frac{v_t}{1-\beta_2^t},\qquad
\theta_t = \theta_{t-1} - \eta\,\frac{\hat m_t}{\sqrt{\hat v_t}+\epsilon}.
\label{eq:optim:adam}
\end{equation}
At $t=1$, $\hat m_1 = g_1$ and $\hat v_1 = g_1^2$, so the first update is
$-\eta\,\mathrm{sign}(g_1)$ in every coordinate whose gradient is much larger than $\epsilon$.
Adam's step size is bounded by roughly $\eta$ per coordinate regardless of the gradient's scale.
That property is the main reason it tolerates the badly scaled gradients of transformers, whose
embedding, attention and MLP gradients differ by orders of magnitude.

\paragraph{Weight decay: coupled versus decoupled.} Classic L2 regularisation adds
$\tfrac{\lambda}{2}\|\theta\|^2$ to the loss, i.e.\ $g_t\leftarrow g_t+\lambda\theta$. For SGD
that is identical to shrinking the weights by $(1-\eta\lambda)$ each step. For Adam it is not:
the decay term passes through the $1/\sqrt{\hat v}$ normalisation, so parameters with large
gradients are barely decayed and parameters with tiny gradients are decayed hard. Loshchilov and
Hutter (2019) proposed \emph{decoupling} the decay from the adaptive update:
\begin{equation}
\theta_t = (1-\eta\lambda)\,\theta_{t-1} - \eta\,\frac{\hat m_t}{\sqrt{\hat v_t}+\epsilon}.
\label{eq:optim:adamw}
\end{equation}
This is AdamW, the default optimizer for LLM pretraining. Note that the decay is multiplied by
the \emph{scheduled} learning rate $\eta$, so it also anneals when the learning rate does; when
you read a config that says ``wd = 0.1'', the per-step shrink factor at peak learning rate
$3\times10^{-4}$ is $3\times10^{-5}$.

\begin{lstlisting}[caption={The AdamW update for one parameter tensor (loop path of \texttt{llmsys.optim.AdamW}). It matches \texttt{torch.optim.AdamW} exactly, including the placement of $\epsilon$ after the bias-corrected square root.},label={lst:optim:adamw}]
st["step"] += 1
t = st["step"]
if wd != 0:
    if group["decoupled"]:
        p.mul_(1 - lr * wd)                    # AdamW: shrink, ignores g
    else:
        g = g.add(p, alpha=wd)                 # Adam + L2: decay via g
m, v = st["exp_avg"], st["exp_avg_sq"]
m.lerp_(g, 1 - b1)                             # m = b1*m + (1-b1)*g
v.mul_(b2).addcmul_(g, g, value=1 - b2)        # v = b2*v + (1-b2)*g*g
denom = (v.sqrt() / math.sqrt(1 - b2 ** t)).add_(eps)
p.addcdiv_(m, denom, value=-lr / (1 - b1 ** t))
\end{lstlisting}

\begin{figure}[t]
\centering
\begin{tikzpicture}
\begin{axis}[width=0.92\linewidth, height=5.4cm, xmin=-4.4, xmax=1.6, ymin=-1.55, ymax=1.75,
  xlabel={$\theta_1$ (shallow direction, $\lambda=1$)}, ylabel={$\theta_2$ (steep, $\lambda=25$)},
  legend style={font=\scriptsize, at={(0.98,0.03)}, anchor=south east}, tick label style={font=\scriptsize},
  label style={font=\small}, axis equal image=false]
\foreach \r in {0.5,1,2,3,4}{
  \addplot[black!25, domain=0:360, samples=80, forget plot] ({\r*cos(x)},{\r/5*sin(x)});}
\addplot[warnred, mark=*, mark size=0.9pt, line width=0.6pt] coordinates {(-4.000,1.500) (-3.700,-1.313) (-3.423,1.148) (-3.166,-1.005) (-2.928,0.879) (-2.709,-0.769) (-2.506,0.673) (-2.318,-0.589) (-2.144,0.515) (-1.983,-0.451) (-1.834,0.395) (-1.697,-0.345) (-1.569,0.302) (-1.452,-0.264) (-1.343,0.231) (-1.242,-0.202) (-1.149,0.177) (-1.063,-0.155) (-0.983,0.136) (-0.909,-0.119) (-0.841,0.104) (-0.778,-0.091) (-0.720,0.079) (-0.666,-0.070) (-0.616,0.061) (-0.570,-0.053)};
\addlegendentry{SGD, $\eta=0.075$}
\addplot[sysblue, mark=square*, mark size=0.9pt, line width=0.6pt] coordinates {(-4.000,1.500) (-3.920,0.750) (-3.778,-0.225) (-3.588,-0.892) (-3.365,-0.980) (-3.119,-0.560) (-2.860,0.056) (-2.595,0.521) (-2.332,0.632) (-2.074,0.405) (-1.827,0.021) (-1.592,-0.297) (-1.373,-0.403) (-1.170,-0.286) (-0.984,-0.050) (-0.816,0.164) (-0.665,0.253) (-0.531,0.198) (-0.413,0.055) (-0.310,-0.087) (-0.222,-0.157) (-0.147,-0.135) (-0.084,-0.049) (-0.032,0.044) (0.010,0.096) (0.044,0.090)};
\addlegendentry{momentum $\mu=0.8$, $\eta=0.02$}
\addplot[keygreen, mark=triangle*, mark size=1.1pt, line width=0.6pt] coordinates {(-4.000,1.500) (-3.600,1.100) (-3.202,0.708) (-2.806,0.333) (-2.415,-0.009) (-2.030,-0.298) (-1.654,-0.514) (-1.289,-0.647) (-0.938,-0.700) (-0.605,-0.684) (-0.293,-0.616) (-0.006,-0.508) (0.254,-0.374) (0.481,-0.225) (0.675,-0.074) (0.833,0.068) (0.955,0.193) (1.041,0.290) (1.092,0.355) (1.111,0.386) (1.099,0.384) (1.061,0.351) (0.999,0.294) (0.918,0.219) (0.820,0.133) (0.710,0.043)};
\addlegendentry{Adam, $\eta=0.4$}
\end{axis}
\end{tikzpicture}
\caption{Twenty-five steps of three optimizers on $f(\theta)=\tfrac12(\theta_1^2+25\theta_2^2)$
starting from $(-4,1.5)$, computed with \code{llmsys.optim}. SGD's learning rate is capped by
the steep direction ($\eta\lambda_{\max}=1.9$), so it zig-zags and crawls along the shallow
one. Momentum accumulates speed along $\theta_1$. Adam normalises each coordinate, so its first
steps move diagonally at speed $\eta$ in both directions; it overshoots and spirals in.}
\label{fig:optim:traj}
\end{figure}

\paragraph{Hyperparameters that matter for LLMs.} Typical pretraining settings are
$\beta_1=0.9$, $\beta_2=0.95$, $\epsilon=10^{-8}$, weight decay $0.1$, and a peak learning rate
between $10^{-4}$ and $10^{-3}$ for 0.1--10B-parameter models (smaller models tolerate larger
rates). Lowering $\beta_2$ from 0.999 to 0.95 shortens the second-moment memory from about 1000
steps to 20. When a rare large gradient arrives, $v$ catches up within a few steps instead of
leaving the effective learning rate wrong for hundreds; in practice this is one of the cheapest
stability fixes for large runs. $\epsilon$ matters only for coordinates whose gradients are
tiny; if you keep $v$ in bf16, $\epsilon=10^{-8}$ underflows relative to $\sqrt{v}$ and
behaves like zero.

% =====================================================================
\section{What an optimizer step costs}\label{sec:optim:cost}\covers{CS336 A1 \S4.3}

The forward and backward passes are compute-bound matrix multiplies. The optimizer step is the
opposite: a long chain of elementwise operations with an arithmetic intensity below one FLOP per
byte. Its time is set entirely by how many bytes it moves.

\begin{napkinmath}[AdamW step time for a 1.1B-parameter model]
Take a Llama-style model with $d=2048$, $L=22$, GQA with 32 query and 4 key/value heads,
$\dff=5632$, $V=32{,}000$ and untied embeddings: $N\approx1.10\times10^9$ parameters
(\code{llmsys.flops.params}). In fp32, one AdamW step reads $p, g, m, v$ (16 bytes per
parameter) and writes $p, m, v$ (12 bytes): \textbf{28 bytes per parameter}. Standard mixed
precision moves the same amount: read bf16 $g$ (2), fp32 master, $m$, $v$ (12); write fp32 master,
$m$, $v$ (12) and the bf16 working copy (2). That is
\[
28 \times 1.10\times10^{9} \approx 30.8\ \text{GB per step}
\;\Rightarrow\; \frac{30.8\ \text{GB}}{3.35\ \text{TB/s}} \approx 9.2\ \text{ms on one H100}.
\]
The arithmetic is about 18 FLOPs per parameter, $2\times10^{10}$ FLOPs in total, which takes 0.02\,ms
at the H100's peak. A training step on a batch of, say, 0.5M tokens costs
$6N\cdot 5\times10^5 = 3.3\times10^{15}$ FLOPs, about 8\,s at 40\% MFU, so the optimizer is
around 0.1\% of the step. If your profile shows a much larger share, the update is not one
memory pass: it is hundreds of small kernels.
\end{napkinmath}

That last case is the usual one for a naive implementation. The loop in \cref{lst:optim:adamw}
launches about ten kernels per parameter tensor, and a transformer has hundreds of tensors (a
28-layer model has around 300). Many are small (norm gains have $d$ elements), so the step can
become \emph{launch-bound}: thousands of kernel launches at a few microseconds each. PyTorch
solves this in two steps:

\begin{itemize}
  \item \textbf{Multi-tensor (``foreach'') kernels} take a \emph{list} of tensors and process
        all of them in one launch, chunking the work across thread blocks. Our
        \code{AdamW(foreach=True)} path (\cref{lst:optim:foreach}) issues about ten launches per
        parameter \emph{group} instead of per tensor, and the test checks it matches the loop to
        floating-point precision.
  \item \textbf{Fused kernels} (\code{torch.optim.AdamW(fused=True)} on CUDA) do the whole update
        for a chunk of parameters in registers: one read and one write of each tensor, the 28
        bytes per parameter of the napkin estimate.
\end{itemize}

\begin{lstlisting}[caption={The same AdamW update over a whole parameter group with multi-tensor kernels.},label={lst:optim:foreach}]
torch._foreach_mul_(params, 1 - lr * wd)                 # decoupled decay
torch._foreach_lerp_(ms, grads, 1 - b1)                  # m
torch._foreach_mul_(vs, b2)
torch._foreach_addcmul_(vs, grads, grads, value=1 - b2)  # v
denoms = torch._foreach_sqrt(vs)
torch._foreach_div_(denoms, [math.sqrt(1 - b2 ** t) for t in steps])
torch._foreach_add_(denoms, eps)
upd = torch._foreach_div(ms, denoms)
torch._foreach_mul_(upd, [-lr / (1 - b1 ** t) for t in steps])
torch._foreach_add_(params, upd)
\end{lstlisting}

\begin{figure}[t]
\centering
\begin{tikzpicture}
\begin{axis}[xbar stacked, width=0.7\linewidth, height=4.8cm, bar width=9pt,
  xlabel={bytes of model state per parameter}, xmin=0, xmax=18.5,
  symbolic y coords={{8-bit Adam states},{bf16 states (pure bf16)},{mixed precision + fp32 master},{fp32 AdamW}},
  ytick=data, y tick label style={font=\scriptsize}, tick label style={font=\scriptsize},
  legend style={font=\scriptsize, at={(0.5,1.06)}, anchor=south, legend columns=5},
  enlarge y limits=0.2]
\addplot[fill=fcompute!60, draw=sysblue] coordinates {(4,{fp32 AdamW}) (2,{mixed precision + fp32 master}) (2,{bf16 states (pure bf16)}) (2,{8-bit Adam states})};
\addplot[fill=fdata!60, draw=keygreen] coordinates {(4,{fp32 AdamW}) (2,{mixed precision + fp32 master}) (2,{bf16 states (pure bf16)}) (2,{8-bit Adam states})};
\addplot[fill=black!25, draw=black!60] coordinates {(0,{fp32 AdamW}) (4,{mixed precision + fp32 master}) (0,{bf16 states (pure bf16)}) (4,{8-bit Adam states})};
\addplot[fill=fmemory!70, draw=napkin] coordinates {(4,{fp32 AdamW}) (4,{mixed precision + fp32 master}) (2,{bf16 states (pure bf16)}) (1,{8-bit Adam states})};
\addplot[fill=fcomm!60, draw=labpurple] coordinates {(4,{fp32 AdamW}) (4,{mixed precision + fp32 master}) (2,{bf16 states (pure bf16)}) (1,{8-bit Adam states})};
\legend{weights, grads, fp32 master, $m$, $v$}
\end{axis}
\end{tikzpicture}
\caption{Bytes of model state per parameter under common precision recipes (activations
excluded). Both fp32 and standard mixed precision cost 16\,B/param; pure-bf16 states halve
that but lose precision in $v$ and in small weight updates; 8-bit optimizers (Dettmers et al.,
2022) quantise $m$ and $v$ block-wise. \Cref{ch:sharding} removes the redundancy across GPUs instead.}
\label{fig:optim:statebytes}
\end{figure}

\begin{systemslens}[Optimizer state is the largest tensor you own]
For a model trained with AdamW, $m$ and $v$ together are twice the size of the fp32 weights, so
the optimizer state is the largest consumer of memory in the job. This drives three design
decisions you will meet later: sharding optimizer state across data-parallel ranks (ZeRO-1,
\cref{ch:sharding}), offloading it to CPU memory, and lower-precision or factored states (8-bit
Adam; Adafactor's row/column factorisation of $v$, \cref{sec:optim:beyond}). It also explains why
checkpoints are about three times the size of the released model: a 7B-parameter model
releases 14\,GB of bf16 weights but writes 112\,GB of training checkpoint.
\end{systemslens}

% =====================================================================
\section{Accounting for memory and compute in a training run}\label{sec:optim:account}\covers{CS336 A1 \S4.3; CS336 L2}

Before launching a run you should be able to predict whether it fits in memory, how long a step
takes, and how many hours the whole run needs. The method has four parts, each a sum over
named tensors or matrix multiplies.

\paragraph{1. Parameters, gradients, optimizer state.} Count parameters by listing every
weight matrix (\cref{sec:gpt:accounting} does it for GPT-2, \cref{ch:modern} for Qwen3). Then
multiply by bytes per parameter for each copy (\cref{fig:optim:statebytes}). For our 1.1B example:
$1.10\times10^9\times16 = 17.6$\,GB of model state.

\paragraph{2. Activations.} Walk through one block and list every tensor that autograd saves
for the backward pass, with its shape. For a pre-norm SwiGLU block with naive attention, on a
batch of $B$ sequences of length $T$:
\begin{center}\small
\begin{tabular}{@{}ll@{}}
\toprule
saved tensor & elements \\
\midrule
inputs to the two RMSNorms & $2BTd$ \\
normalised input to $W_Q,W_K,W_V$ & $BTd$ \\
$Q$, $K$, $V$ (GQA) & $BT(H+2G)\dhead$ \\
attention scores and probabilities & $2BHT^2$ \\
attention output (input to $W_O$) & $BTH\dhead$ \\
normalised input to the MLP & $BTd$ \\
gate and up projections; SiLU output; their product & $4BT\dff$ \\
\bottomrule
\end{tabular}
\end{center}
Add the final norm, the logits ($BTV$) and the softmax needed by the loss's backward. Which
tensors are actually saved depends on the implementation: fused kernels and FlashAttention save
fewer. \Cref{ch:perf} shows how to measure them with saved-tensor hooks.

\paragraph{3. FLOPs.} $2mnp$ for every $(m\times n)(n\times p)$ product, summed over the forward
pass; the backward pass costs about twice the forward, giving the familiar $6N$ per token plus
the attention term $\approx 12 L H \dhead T$ per token (forward + backward) that grows with
context. The optimizer adds $\approx18N$ FLOPs per step, independent of batch size.

\paragraph{4. Time.} Divide total FLOPs by achievable throughput, i.e.\ peak FLOP/s times the
model-FLOPs utilisation (MFU) you expect (35--50\% is good for dense transformers).

\begin{napkinmath}[Fitting and timing the 1.1B example]
With $T=2048$ in bf16 and naive attention, the table gives about $3.4\times10^8$ elements, or
0.68\,GB, per layer per sequence, and 15.5\,GB per sequence for all 22 layers plus logits. The
$2BHT^2$ attention matrices are 79\% of that. With FlashAttention (which saves only the output
and a log-sum-exp per row) it falls to about 0.145\,GB per layer, 3.7\,GB per sequence. On an
80\,GB GPU with 17.6\,GB of model state, the activation budget allows about 16 sequences of 2048
tokens per micro-batch with FlashAttention, but only 4 with naive attention. Either way the
peak memory is linear in batch size, $\text{mem}(B) = a\cdot B + b$ with $b$ the model state, which is the form you should always
derive.

Compute: $6N + 12LH\dhead T \approx 6.6 + 0.7 = 7.3$\,GFLOPs per token at $T=2048$. Training on
$10^{11}$ tokens costs $\approx 7.3\times10^{20}$ FLOPs; on 8 H100s at 40\% MFU
($8\times 989\times10^{12}\times0.4 = 3.2\times10^{15}$ FLOP/s) that is about 64 hours.
\end{napkinmath}

\begin{pitfall}[``Float32'' peak numbers]
Vendor tables often quote a tensor-core number for ``FP32'' that is really TF32: 10-bit
mantissa inputs with fp32 accumulation. On an H100 SXM that figure is about half the bf16 dense
peak, and true IEEE fp32 matrix multiplies run at a much lower rate. When you compute MFU, use
the peak that matches the arithmetic you actually run and say which one it is. With
\code{torch.set\_float32\_matmul\_precision("high")} PyTorch uses TF32 for fp32 matmuls on Ampere
and later GPUs.
\end{pitfall}

% =====================================================================
\section{Learning-rate schedules}\label{sec:optim:sched}\covers{CS336 A1 \S4.4}

A schedule is a function $\eta(t)$ evaluated every step and written into
\code{group["lr"]}. Three shapes cover almost all LLM training (\cref{fig:optim:sched}).

\begin{description}[leftmargin=1.2em, style=nextline]
  \item[Linear warmup.] Ramp from near zero to $\eta_{\max}$ over $T_w$ steps (typically
    0.1--2\% of training, often 1000--2000 steps). Adam's second-moment estimate is noisy for
    the first $\sim 1/(1-\beta_2)$ steps, and at initialisation the residual stream and attention
    logits are far from their trained scales. Large early steps at full learning rate can push
    the model into a regime it never recovers from (loss spikes, dead attention heads).
  \item[Cosine decay.] $\eta(t) = \eta_{\min} + \tfrac12(\eta_{\max}-\eta_{\min})
    \bigl(1+\cos(\pi\frac{t-T_w}{T_c-T_w})\bigr)$ for $T_w\le t\le T_c$, then constant
    $\eta_{\min}$. It is the default since GPT-3 and Chinchilla. The run length $T_c$ is baked in:
    if you stop early you get a model trained at a high learning rate, and if you continue past
    $T_c$ you train at $\eta_{\min}$ forever. When you change the step budget (for example to run
    on a smaller machine), set $T_c$ to the new final step.
  \item[Warmup--stable--decay (WSD).] Hold $\eta_{\max}$ constant for most of training and decay
    over the last 10--20\%, linearly or as $1-\sqrt{x}$ (Hu et al., 2024; H\"agele et al., 2024).
    Because the long stable phase does not depend on the final length, you can branch decay
    phases off a single run at several token budgets. \Cref{ch:scalinglaws} uses this to make
    scaling-law sweeps several times cheaper.
\end{description}

\begin{figure}[t]
\centering
\begin{tikzpicture}
\begin{axis}[width=0.92\linewidth, height=4.8cm, xmin=0, xmax=1000, ymin=0, ymax=1.08,
  xlabel={step}, ylabel={$\eta(t)/\eta_{\max}$}, tick label style={font=\scriptsize},
  label style={font=\small}, legend style={font=\scriptsize, at={(0.02,0.03)}, anchor=south west}]
\addplot[sysblue, line width=0.9pt, domain=0:1000, samples=201]
  {x<50 ? (x+1)/50 : 0.1 + 0.45*(1+cos(deg(pi*(x-50)/950)))};
\addlegendentry{warmup + cosine (floor 0.1)}
\addplot[keygreen, line width=0.9pt, domain=0:1000, samples=201]
  {x<50 ? (x+1)/50 : (x<800 ? 1 : 0.1 + 0.9*(1-(x-800)/200))};
\addlegendentry{WSD, linear decay over last 20\%}
\addplot[napkin, line width=0.9pt, dashed, domain=0:1000, samples=201]
  {x<50 ? (x+1)/50 : (x<800 ? 1 : 0.1 + 0.9*(1-sqrt((x-800)/200)))};
\addlegendentry{WSD, $1-\sqrt{x}$ decay}
\addplot[warnred, line width=0.9pt, densely dotted, domain=0:1000, samples=201]
  {x<50 ? (x+1)/50 : 1 - 0.9*(x-50)/950};
\addlegendentry{warmup + linear decay}
\end{axis}
\end{tikzpicture}
\caption{The schedules implemented in \code{llmsys.optim}, for a 1000-step run with 50 warmup
steps and a floor of $0.1\,\eta_{\max}$.}
\label{fig:optim:sched}
\end{figure}

\begin{lstlisting}[caption={Schedules are plain callables, so they are trivial to test, plot and checkpoint.},label={lst:optim:sched}]
@dataclass
class WarmupCosine:
    peak: float; total: int; warmup: int = 0; floor: float = 0.0
    def __call__(self, step):
        if self.warmup > 0 and step < self.warmup:
            return self.peak * (step + 1) / self.warmup      # step 0 already moves
        if step >= self.total:
            return self.floor
        prog = (step - self.warmup) / max(1, self.total - self.warmup)
        return self.floor + 0.5 * (self.peak - self.floor) * (1 + math.cos(math.pi * prog))

sched = WarmupCosine(peak=3e-4, total=20_000, warmup=500, floor=3e-5)
for step in range(start, 20_000):
    set_lr(opt, sched(step))      # writes every param group (times an optional per-group scale)
\end{lstlisting}

Two conventions differ between codebases and cause off-by-one disagreements. Is the first step
0 or 1? Does warmup reach $\eta_{\max}$ at step $T_w-1$ or $T_w$? Ours is 0-based and reaches the peak at
step $T_w-1$. Any convention is fine if your tests pin it down, and if a resumed run calls
the schedule with the restored step rather than a step counter that restarted at zero.

We measured the WSD branching idea on CPU with Lab~9's \code{run.py wsd} command
(\cref{fig:optim:wsd}).

\begin{figure}[t]
\centering
\begin{tikzpicture}
\begin{axis}[width=0.9\linewidth, height=5.2cm, xmin=0, xmax=610, ymin=1.45, ymax=1.85,
  xlabel={step}, ylabel={val loss (nats/token)}, tick label style={font=\scriptsize},
  label style={font=\small}, legend style={font=\scriptsize, at={(0.02,0.03)}, anchor=south west}]
\addplot[black!45, line width=0.7pt] coordinates {(99,1.81) (119,1.727) (139,1.681) (159,1.654) (179,1.622) (199,1.619) (219,1.601) (239,1.591) (259,1.591) (279,1.583) (299,1.58) (319,1.579) (339,1.574) (359,1.58) (379,1.57) (399,1.558) (419,1.562) (439,1.56) (459,1.556) (479,1.564)};
\addlegendentry{WSD trunk (constant $\eta$)}
\addplot[keygreen, line width=1.0pt, forget plot] coordinates {(159,1.654) (179,1.613) (199,1.579)};
\addplot[keygreen, line width=1.0pt, forget plot] coordinates {(319,1.579) (339,1.56) (359,1.54) (379,1.524) (399,1.505)};
\addplot[keygreen, line width=1.0pt] coordinates {(479,1.564) (499,1.551) (519,1.539) (539,1.514) (559,1.508) (579,1.496) (599,1.488)};
\addlegendentry{WSD decay branches (20\%)}
\addplot[sysblue, line width=0.7pt, densely dashed] coordinates {(99,1.812) (119,1.727) (139,1.68) (159,1.651) (179,1.622) (199,1.616) (219,1.588) (239,1.578) (259,1.574) (279,1.558) (299,1.549) (319,1.55) (339,1.54) (359,1.528) (379,1.527) (399,1.516) (419,1.511) (439,1.514) (459,1.507) (479,1.502) (499,1.501) (519,1.495) (539,1.489) (559,1.49) (579,1.488) (599,1.486)};
\addlegendentry{cosine run, 600-step budget}
\addplot[only marks, mark=*, sysblue, mark size=1.8pt] coordinates {(199,1.638) (399,1.512) (599,1.486)};
\addlegendentry{cosine runs' finals (200/400/600)}
\addplot[only marks, mark=square*, keygreen, mark size=1.8pt, forget plot] coordinates {(199,1.579) (399,1.505) (599,1.488)};
\end{axis}
\end{tikzpicture}
\caption{One WSD trunk with three decay branches versus three separate cosine runs, at budgets
of 200, 400 and 600 steps (same model and data as \cref{fig:optim:lrsweep}, $\eta=3\times10^{-3}$;
measured on CPU, one seed). The branches end at 1.579, 1.505 and 1.488, against 1.638, 1.512 and
1.486 for the cosine runs, so they match or beat them. They cost 720 optimizer steps in total
(480 trunk + 40 + 80 + 120 decay) instead of 1200. The loss falls steeply as each branch decays:
much of the benefit usually credited to a schedule's shape comes from annealing at the end.}
\label{fig:optim:wsd}
\end{figure}

% =====================================================================
\section{Gradient clipping}\label{sec:optim:clip}\covers{CS336 A1 \S4.5}

Occasionally a batch produces a gradient far larger than usual: a rare token, a pathological
document, or the model entering a sharp region. One such step can undo hours of progress.
Global-norm clipping rescales the \emph{entire} gradient vector when its norm exceeds $M$:
\begin{equation}
\vect{g} \leftarrow \vect{g}\cdot\min\Bigl(1, \frac{M}{\|\vect{g}\|_2+\epsilon}\Bigr),
\qquad \|\vect{g}\|_2 = \Bigl(\sum_{\text{tensors }i}\|\vect{g}_i\|_2^2\Bigr)^{1/2}.
\label{eq:optim:clip}
\end{equation}
Clipping the global norm preserves the gradient's \emph{direction}; clipping each tensor (or
each element) separately would change it (\cref{fig:optim:clip}). $M=1.0$ is the near-universal
LLM default. With Adam, clipping matters less for the direction (Adam normalises per coordinate
anyway) than for the moments: a huge $g$ would inflate $v$ for the next $\sim 1/(1-\beta_2)$
steps.

\begin{figure}[t]
\centering
\begin{tikzpicture}[scale=1.25]
\draw[->, black!50] (-0.3,0)--(3.4,0) node[right,font=\scriptsize]{$g_1$};
\draw[->, black!50] (0,-0.3)--(0,2.6) node[above,font=\scriptsize]{$g_2$};
\draw[sysblue, line width=0.6pt, fill=sysbluebg] (0,0) circle (1);
\node[font=\scriptsize, sysblue] at (-0.55,-0.6) {$\|g\|\le M$};
\draw[-{Stealth}, warnred, line width=1pt] (0,0)--(3.0,1.8) node[right,font=\scriptsize]{$g$, $\|g\|=3.5M$};
\draw[-{Stealth}, keygreen, line width=1.2pt] (0,0)--(0.857,0.514) node[below right,font=\scriptsize, xshift=-2pt]{global clip};
\draw[-{Stealth}, napkin, line width=1pt, dashed] (0,0)--(1,1) node[above,font=\scriptsize]{per-coordinate clip};
\draw[black!40, dotted] (1,0)--(1,1)--(0,1);
\end{tikzpicture}
\caption{Global-norm clipping (green) shrinks the gradient onto the ball of radius $M$ and keeps
its direction. Clipping each coordinate (or each tensor) to $[-M,M]$ (dashed) changes the
direction, here from $31^\circ$ to $45^\circ$.}
\label{fig:optim:clip}
\end{figure}

\begin{lstlisting}[caption={Global-norm clipping with multi-tensor kernels. It returns the \emph{pre-clip} norm, the single most useful scalar to log.},label={lst:optim:clipcode}]
@torch.no_grad()
def clip_grad_norm_(parameters, max_norm, eps=1e-6, foreach=True):
    grads = [p.grad for p in parameters if p.grad is not None]
    norms = torch._foreach_norm([g.float() for g in grads])   # fp32 norms
    total = torch.linalg.vector_norm(torch.stack(norms))      # global ||g||
    scale = max_norm / (total + eps)
    if scale < 1:                     # NOTE: syncs GPU -> CPU
        torch._foreach_mul_(grads, scale.to(grads[0].dtype))
    return float(total)
\end{lstlisting}

\begin{systemslens}[Clipping is a global reduction]
\Cref{eq:optim:clip} needs a norm over \emph{all} parameters before any of them can be scaled.
On one GPU that is one reduction kernel, but turning the result into a Python \code{bool} or
\code{float} forces a host--device synchronisation. PyTorch's own implementation avoids the
branch by always multiplying by \code{clamp(scale, max=1)} on the device. When parameters are
sharded across GPUs (FSDP, tensor parallelism, \cref{ch:sharding,ch:distributed}), each rank
holds only part of $\vect{g}$, so clipping needs an all-reduce of the squared partial norms.
That is one more collective per step, and the place where many custom parallel trainers get the
semantics subtly wrong by clipping each shard separately. With fp16 loss scaling
(\cref{sec:pretrain:mixed}), always \emph{unscale} gradients before computing the norm.
\end{systemslens}

% =====================================================================
\section{Beyond AdamW: Lion, Adafactor, Shampoo and Muon}\label{sec:optim:beyond}

AdamW has been the default for LLMs since about 2019, but it is not the only option, and recent
speed-running efforts on small transformers have favoured alternatives. A short survey, from
simplest to most involved:

\begin{description}[leftmargin=1.2em, style=nextline]
  \item[Lion] (Chen et al., 2023) was found by program search: take the \emph{sign} of an
    interpolation between momentum and gradient, $\theta \leftarrow \theta -
    \eta\,\mathrm{sign}(\beta_1 m + (1-\beta_1) g)$, then update $m$ with a second rate
    $\beta_2$. It needs one state tensor instead of two (12 vs 16 bytes/param in mixed precision) and
    a learning rate 3--10$\times$ smaller than AdamW's, because every coordinate moves by exactly
    $\eta$.
  \item[Adafactor] (Shazeer and Stern, 2018) stores the second moment of an $n\times m$ matrix
    as a rank-1 factorisation (row and column sums, $n+m$ numbers instead of $nm$), with optional
    update clipping and relative step sizes. It was the memory-saving default for T5 and PaLM.
  \item[Shampoo and SOAP] keep a left and a right preconditioner per weight matrix,
    $\theta\leftarrow\theta-\eta\,L^{-1/4} G R^{-1/4}$ with $L=\sum GG^\top$ and
    $R=\sum G^\top G$. This approximates full-matrix AdaGrad for the matrix's Kronecker
    structure. The inverse roots are expensive and are recomputed every few steps, often on
    separate hardware; SOAP (Vyas et al., 2024) runs Adam in Shampoo's eigenbasis.
  \item[Muon] (Jordan et al., 2024) keeps a momentum buffer for each hidden weight matrix and
    replaces the update by its nearest semi-orthogonal matrix: if the momentum is $U\Sigma V^\top$,
    the update is $UV^\top$, which equalises all singular directions. The orthogonalisation uses a few
    Newton--Schulz iterations of an odd polynomial (\cref{fig:optim:ns}), five matrix multiplies on the smaller
    Gram matrix, so it runs in bf16 on tensor cores. Embeddings, the LM head and all 1-D
    parameters stay on AdamW.
\end{description}

\begin{figure}[t]
\centering
\begin{tikzpicture}
\begin{axis}[width=0.86\linewidth, height=4.8cm, xmin=0, xmax=1, ymin=0, ymax=1.32,
  xlabel={input singular value $\sigma$ (after Frobenius normalisation)}, ylabel={output $\sigma'$},
  tick label style={font=\scriptsize}, label style={font=\small},
  legend style={font=\scriptsize, at={(0.98,0.03)}, anchor=south east}]
\addplot[fill=keygreen!12, draw=none, forget plot] coordinates {(0,0.68) (1,0.68) (1,1.2) (0,1.2)} \closedcycle;
\addplot[black!40, dashed, forget plot] coordinates {(0,1) (1,1)};
\addplot[black!45, line width=0.6pt] coordinates {(0.000,0.000) (0.005,0.017) (0.010,0.034) (0.015,0.052) (0.020,0.069) (0.025,0.086) (0.030,0.103) (0.035,0.120) (0.040,0.137) (0.045,0.155) (0.050,0.172) (0.060,0.206) (0.080,0.273) (0.100,0.340) (0.120,0.405) (0.140,0.469) (0.160,0.532) (0.180,0.593) (0.200,0.651) (0.220,0.708) (0.240,0.762) (0.260,0.814) (0.280,0.863) (0.300,0.909) (0.320,0.953) (0.340,0.993) (0.360,1.030) (0.380,1.063) (0.400,1.093) (0.420,1.119) (0.440,1.142) (0.460,1.162) (0.480,1.177) (0.500,1.189) (0.520,1.197) (0.540,1.201) (0.560,1.202) (0.580,1.199) (0.600,1.193) (0.620,1.184) (0.640,1.171) (0.660,1.155) (0.680,1.136) (0.700,1.115) (0.720,1.091) (0.740,1.065) (0.760,1.037) (0.780,1.007) (0.800,0.976) (0.820,0.945) (0.840,0.913) (0.860,0.881) (0.880,0.849) (0.900,0.819) (0.920,0.790) (0.940,0.763) (0.960,0.739) (0.980,0.718) (1.000,0.701)};
\addlegendentry{1 iteration}
\addplot[sysblue, line width=0.7pt] coordinates {(0.000,0.000) (0.005,0.059) (0.010,0.118) (0.015,0.177) (0.020,0.236) (0.025,0.293) (0.030,0.350) (0.035,0.406) (0.040,0.461) (0.045,0.515) (0.050,0.567) (0.060,0.668) (0.080,0.847) (0.100,0.992) (0.120,1.100) (0.140,1.169) (0.160,1.200) (0.180,1.196) (0.200,1.162) (0.220,1.105) (0.240,1.033) (0.260,0.954) (0.280,0.876) (0.300,0.805) (0.320,0.747) (0.340,0.707) (0.360,0.685) (0.380,0.683) (0.400,0.699) (0.420,0.729) (0.440,0.769) (0.460,0.813) (0.480,0.857) (0.500,0.896) (0.520,0.926) (0.540,0.943) (0.560,0.946) (0.580,0.935) (0.600,0.912) (0.620,0.879) (0.640,0.839) (0.660,0.797) (0.680,0.757) (0.700,0.722) (0.720,0.697) (0.740,0.684) (0.760,0.683) (0.780,0.696) (0.800,0.721) (0.820,0.757) (0.840,0.800) (0.860,0.848) (0.880,0.898) (0.900,0.947) (0.920,0.993) (0.940,1.033) (0.960,1.067) (0.980,1.094) (1.000,1.114)};
\addlegendentry{2 iterations}
\addplot[warnred, line width=0.7pt] coordinates {(0.000,0.000) (0.005,1.154) (0.010,0.699) (0.015,0.931) (0.020,0.682) (0.025,0.933) (0.030,1.122) (0.035,1.107) (0.040,0.951) (0.045,0.791) (0.050,0.758) (0.060,1.006) (0.080,0.952) (0.100,0.712) (0.120,0.715) (0.140,0.786) (0.160,1.028) (0.180,0.997) (0.200,0.747) (0.220,0.707) (0.240,0.749) (0.260,0.682) (0.280,0.838) (0.300,1.080) (0.320,1.134) (0.340,1.100) (0.360,1.056) (0.380,1.050) (0.400,1.086) (0.420,1.126) (0.440,1.129) (0.460,1.060) (0.480,0.911) (0.500,0.765) (0.520,0.698) (0.540,0.683) (0.560,0.682) (0.580,0.687) (0.600,0.723) (0.620,0.827) (0.640,0.981) (0.660,1.096) (0.680,1.134) (0.700,1.120) (0.720,1.082) (0.740,1.052) (0.760,1.051) (0.780,1.080) (0.800,1.119) (0.820,1.134) (0.840,1.090) (0.860,0.947) (0.880,0.760) (0.900,0.682) (0.920,0.713) (0.940,0.749) (0.960,0.749) (0.980,0.723) (1.000,0.696)};
\addlegendentry{5 iterations}
\end{axis}
\end{tikzpicture}
\caption{The Newton--Schulz polynomial $p(\sigma)=a\sigma+b\sigma^3+c\sigma^5$ with Muon's
coefficients $(3.4445,-4.7750,2.0315)$, iterated. It acts on each singular value independently
because $X\mapsto aX+b(XX^\top)X+c(XX^\top)^2X$ is an odd matrix polynomial. After five
iterations every singular value above $\sim$0.005 lands in the shaded band $[0.68,1.2]$: not
exactly 1, but close enough, and the coefficients were tuned for speed of reaching the band
rather than exact convergence. Our test measures $[0.68,1.13]$ on a random $32\times16$ matrix.}
\label{fig:optim:ns}
\end{figure}

\begin{lstlisting}[caption={Newton--Schulz orthogonalisation, the core of Muon.},label={lst:optim:ns}]
def newton_schulz_orthogonalize(G, steps=5, eps=1e-7):
    a, b, c = 3.4445, -4.7750, 2.0315
    X = G.float()
    if X.shape[0] > X.shape[1]:            # work with the smaller Gram matrix
        X = X.T
    X = X / (X.norm() + eps)               # ||X||_F >= sigma_max, so all sigma <= 1
    for _ in range(steps):
        A = X @ X.T                        # (r, r), r = min(m, n)
        X = a * X + (b * A + c * A @ A) @ X
    return (X.T if G.shape[0] > G.shape[1] else X).to(G.dtype)
\end{lstlisting}

\begin{napkinmath}[What Muon's orthogonalisation costs]
For a $d\times d$ weight, one iteration costs $X X^\top$ ($2d^3$), $A^2$ ($2d^3$) and the product
with $X$ ($2d^3$), so $6d^3$ FLOPs, or $30d^3$ for five iterations. Per optimizer step, over all hidden matrices of a
model with $N$ parameters arranged as square $d\times d$ blocks, that is $30 d\cdot N$ FLOPs.
For $d=2048$ that is $6\times10^4 N$, the same as the forward and backward pass over
$10^4$ tokens. With a batch of $5\times10^5$ tokens it adds about 2\% to the step, and it is
compute-bound tensor-core work, unlike AdamW's memory-bound update. The real systems cost
is elsewhere: the update needs the \emph{whole} matrix, so under FSDP or tensor parallelism
(\cref{ch:sharding}) each matrix has to be gathered or the orthogonalisation distributed.
\end{napkinmath}

\Cref{fig:optim:compare} compares the four optimizers on a small task we can train on a CPU.
Take it as a demonstration that the implementations work, not as a ranking: each optimizer
received one hand-picked learning rate, and the model has two layers.

\begin{figure}[t]
\centering
\begin{tikzpicture}
\begin{axis}[width=0.9\linewidth, height=5.0cm, xmin=0, xmax=400, ymin=1.4, ymax=3.0,
  xlabel={step}, ylabel={val loss (nats/token)}, tick label style={font=\scriptsize},
  label style={font=\small}, legend style={font=\scriptsize, at={(0.98,0.98)}, anchor=north east}]
\addplot[black!40, dashed, forget plot] coordinates {(0,1.443) (400,1.443)};
\addplot[warnred, line width=0.8pt] coordinates {(24,3.394) (49,1.989) (74,1.734) (99,1.679) (124,1.634) (149,1.607) (174,1.59) (199,1.572) (224,1.556) (249,1.538) (274,1.524) (299,1.511) (324,1.505) (349,1.49) (374,1.489) (399,1.484)};
\addlegendentry{SGD + Nesterov, $\eta=0.3$ (final 1.484)}
\addplot[sysblue, line width=0.8pt] coordinates {(24,2.731) (49,1.898) (74,1.745) (99,1.744) (124,1.688) (149,1.659) (174,1.639) (199,1.618) (224,1.599) (249,1.569) (274,1.545) (299,1.528) (324,1.519) (349,1.499) (374,1.498) (399,1.49)};
\addlegendentry{AdamW, $\eta=10^{-2}$ (1.490)}
\addplot[keygreen, line width=0.8pt] coordinates {(24,3.784) (49,2.388) (74,1.947) (99,1.82) (124,1.727) (149,1.681) (174,1.651) (199,1.619) (224,1.584) (249,1.566) (274,1.55) (299,1.536) (324,1.528) (349,1.515) (374,1.508) (399,1.504)};
\addlegendentry{Muon $\eta=0.05$ + AdamW (1.504)}
\addplot[napkin, line width=0.8pt] coordinates {(24,5.156) (49,3.864) (74,3.025) (99,2.486) (124,2.146) (149,1.931) (174,1.797) (199,1.715) (224,1.658) (249,1.618) (274,1.595) (299,1.574) (324,1.563) (349,1.553) (374,1.545) (399,1.538)};
\addlegendentry{Lion, $\eta=10^{-3}$ (1.538)}
\end{axis}
\end{tikzpicture}
\caption{Four optimizers from \code{llmsys.optim} on the task of \cref{fig:optim:lrsweep},
each with one learning rate picked by hand (measured on CPU, one seed). All four reach within
0.06 nats of each other, and plain SGD with Nesterov momentum comes out ahead. On a two-layer
model with a stationary data distribution, the gradient-scale problems that motivate adaptive
methods barely arise. Do not read rankings from toy runs; the lab asks you to tune each optimizer
before comparing.}
\label{fig:optim:compare}
\end{figure}


% =====================================================================
\section{Feeding the loop: batches from a memory-mapped token array}\label{sec:optim:data}\covers{CS336 A1 \S5.1}

Pretraining data is tokenised once, offline, into one long array of token ids with an
end-of-text separator between documents, and stored as raw little-endian integers
(\cref{sec:text:pipeline}). With a vocabulary under 65,536, \code{uint16} halves the file size
relative to \code{int32}; a 10B-token corpus is 20\,GB. The training loop then samples random
windows from it (\cref{fig:optim:window}):

\begin{lstlisting}[caption={Random windows from a memory-mapped token file (\texttt{RandomWindowBatcher.\_\_next\_\_}).},label={lst:optim:batcher}]
tokens = np.memmap("train.bin", dtype=np.uint16, mode="r")   # nothing read yet
starts = rng.integers(0, len(tokens) - T, size=B)            # B random offsets
idx = starts[:, None] + np.arange(T + 1)[None, :]            # (B, T+1) indices
chunk = torch.from_numpy(tokens[idx].astype(np.int64))       # reads ~B*(T+1)*2 bytes
chunk = chunk.pin_memory().to("cuda", non_blocking=True)     # async H2D copy
x, y = chunk[:, :-1], chunk[:, 1:]                           # targets = inputs shifted by 1
\end{lstlisting}

\begin{figure}[t]
\centering
\begin{tikzpicture}[x=0.38cm, y=0.5cm]
\foreach \i in {0,...,33} {\draw[black!40, fill=fdata!10] (\i,0) rectangle +(1,1);}
\foreach \i in {7,21} {\fill[warnred!35] (\i,0) rectangle +(1,1);}
\node[font=\scriptsize, anchor=east] at (-0.2,0.5) {\code{train.bin}};
\node[font=\tiny, warnred] at (7.5,1.45) {EOT}; \node[font=\tiny, warnred] at (21.5,1.45) {EOT};
\draw[decorate, decoration={brace, amplitude=4pt, mirror}, sysblue, line width=0.6pt] (3,-0.15)--(11,-0.15);
\draw[decorate, decoration={brace, amplitude=4pt, mirror}, keygreen, line width=0.6pt] (18,-0.15)--(26,-0.15);
\draw[sysblue, line width=0.8pt] (3,0) rectangle (11,1);
\draw[keygreen, line width=0.8pt] (18,0) rectangle (26,1);
\node[font=\scriptsize, sysblue] at (7,-1.0) {start $s_1$, $T{+}1$ tokens};
\node[font=\scriptsize, keygreen] at (22,-1.0) {start $s_2$};
\begin{scope}[shift={(4,-4.2)}]
\foreach \i in {0,...,7} {\draw[sysblue, fill=sysbluebg] (\i,0.6) rectangle +(1,1);}
\foreach \i in {1,...,8} {\draw[napkin, fill=napkinbg] (\i,-0.9) rectangle +(1,1);}
\node[font=\scriptsize, anchor=east] at (-0.2,1.1) {$x$};
\node[font=\scriptsize, anchor=east] at (0.8,-0.4) {$y$};
\node[font=\scriptsize, anchor=west, text width=6.2cm] at (10,0.3) {$y_t = x_{t+1}$: every position is a training example. A window may cross an EOT; the model learns that EOT resets context.};
\end{scope}
\draw[-{Stealth}, black!50] (7,-1.4)--(8,-2.5);
\end{tikzpicture}
\caption{Sampling windows from a flat token file. Each batch row is $T+1$ contiguous tokens
starting at a uniformly random offset; inputs and targets are the two shifted views. The memmap
reads only the touched pages.}
\label{fig:optim:window}
\end{figure}

Why random windows rather than the epoch-based sliding-window dataset of \cref{ch:text}? At
pretraining scale you rarely see a token more than once, so ``epochs'' lose their meaning. Sampling with
replacement is simple, needs no shuffle buffer, and makes every offset equally likely. It also
makes the sampler's state tiny: a random generator. Three engineering details matter:

\begin{itemize}
  \item \textbf{Validate the memmap.} A memmap with the wrong \code{dtype} or a header you forgot
        about silently yields garbage. Assert that \code{tokens.max() < vocab\_size} on a sample and decode
        a few windows back to text before training.
  \item \textbf{Make the sampler resumable.} Use a private generator
        (\code{np.random.default\_rng(seed)}) and save its \code{bit\_generator.state} in the
        checkpoint. A global \code{np.random} call elsewhere in your code would otherwise change
        which batches a resumed run sees.
  \item \textbf{Overlap host work with the GPU.} Fancy indexing a memmap, converting to
        \code{int64} and copying to the device takes milliseconds of host time per batch. Pinned
        memory plus \code{non\_blocking=True} makes the copy asynchronous, and a background thread
        (\code{PrefetchBatcher}) samples the next batch while the GPU computes the current one.
\end{itemize}

\begin{napkinmath}[Does the sampler keep up?]
A batch of $B=64$ windows of $T=1024$ \code{uint16} tokens is 131\,KB read from the page cache
(or a fast NVMe drive). The model-side work for that batch on a 125M-parameter model is
$6\times1.25\times10^8\times65{,}536 \approx 4.9\times10^{13}$ FLOPs, about 0.1\,s on an H100 at
50\% MFU. The data side is three orders of magnitude cheaper, so a single prefetch thread is
plenty. The picture changes only for tiny models or when you tokenise on the fly; tokenise
offline.
\end{napkinmath}

% =====================================================================
\section{Checkpointing everything}\label{sec:optim:ckpt}\covers{CS336 A1 \S5.2}

A checkpoint is correct when \emph{loading it and continuing produces the same run as never
stopping}. That requires more than the weights (\cref{fig:optim:ckpt}):

\begin{figure}[t]
\centering
\begin{tikzpicture}[
  part/.style={draw, rounded corners=2pt, minimum width=3.3cm, minimum height=0.55cm, font=\scriptsize, align=left, text width=3.1cm},
  arr/.style={-{Stealth[length=2mm]}, line width=0.6pt}]
\node[part, fill=fcompute!18] (w) at (0,0) {\code{model}: weights};
\node[part, fill=fmemory!20, below=2pt of w] (o) {\code{optimizer}: $m,v$, step $t$};
\node[part, fill=fdata!18, below=2pt of o] (d) {\code{data}: sampler RNG, count};
\node[part, fill=fcomm!16, below=2pt of d] (r) {\code{rng}: python/numpy/torch/cuda};
\node[part, fill=black!6, below=2pt of r] (e) {\code{step}, \code{extra}: config, best val};
\node[draw, dashed, rounded corners, fit=(w)(e), inner sep=4pt, label={[font=\scriptsize]above:checkpoint dict}] (ck) {};
\node[draw, rounded corners=2pt, font=\scriptsize, align=center, fill=white] (tmp) at (5.6,-0.4) {\code{ckpt\_00004000.pt.tmp}};
\node[draw, rounded corners=2pt, font=\scriptsize, align=center, fill=keygreen!15] (fin) at (5.6,-1.7) {\code{ckpt\_00004000.pt}};
\node[draw, rounded corners=2pt, font=\scriptsize, align=center, fill=black!8, text=black!50] (old) at (5.6,-3.0) {\code{ckpt\_00002000.pt} (deleted)};
\draw[arr] (ck.east |- tmp) -- node[above,font=\tiny]{1. torch.save} (tmp);
\draw[arr] (tmp) -- node[right,font=\tiny]{2. os.replace (atomic)} (fin);
\draw[arr, dashed] (fin) -- node[right,font=\tiny]{3. keep last $k$} (old);
\end{tikzpicture}
\caption{Anatomy of a resumable checkpoint and the write protocol. Writing to a temporary
name and renaming means a crash during step 1 leaves the previous checkpoint intact; readers
never see a half-written file.}
\label{fig:optim:ckpt}
\end{figure}

\begin{description}[leftmargin=1.2em, style=nextline]
  \item[Optimizer state.] AdamW's $m$, $v$ and step count. Restarting with fresh moments re-runs
    the warmup transient: Adam's first steps are sign-like (\cref{sec:optim:adam}) and too
    large for a trained model, and the bias-correction counter also restarts.
  \item[The step and schedule position.] The learning rate is a function of the global step.
  \item[Data position.] For our sampler, its generator state; for a sequential loader, the byte
    offset and shuffle seed. Without it the resumed run revisits some data and skips other data.
  \item[RNG states.] Dropout masks and any sampling inside the loop draw from torch's
    generators. For bit-exactness, restore python, numpy, torch-CPU and every CUDA device's
    generator.
\end{description}

\begin{lstlisting}[caption={Saving and restoring with \texttt{CheckpointManager}.},label={lst:optim:ckptcode}]
ckpt = CheckpointManager("runs/ts-small/ckpt", keep_last=3)
if ckpt.latest() is not None:                       # resume if possible
    start = ckpt.restore(model, opt, data)["step"]
...
if (step + 1) % ckpt_every == 0:                    # step+1 = number of completed steps
    ckpt.save(step + 1, model, opt, data, extra=dict(config=asdict(cfg)))

# inside save(): write-then-rename, so a crash never corrupts the last good file
torch.save(obj, final + ".tmp")
os.replace(final + ".tmp", final)
\end{lstlisting}

The lab's acceptance test is the strongest specification of ``correct'': train 16 steps
straight through; separately, train with a simulated crash at step 11, rebuild the model from
scratch in a ``new process'', resume from the step-8 checkpoint, and require the resumed
losses and final weights to be \emph{bit-identical} to the uninterrupted run. That test catches
every item in the list above, and also loaders that are not seeded.

\begin{systemslens}[Checkpoint I/O at scale]
A checkpoint of a model with $N$ parameters trained with AdamW is about $16N$ bytes: 1.1\,TB
for a 70B model. Writing it from one process at 2\,GB/s takes nine minutes, during which the
GPUs idle. Large systems therefore (i) shard the write so every rank saves its own slice in parallel
(PyTorch Distributed Checkpoint), (ii) copy state to host memory and write asynchronously
while training continues, and (iii) save more often to fast local disk and less often to remote
object storage. The checkpoint interval trades lost work against I/O overhead: with mean time
between failures $M$ and checkpoint cost $C$, Young's approximation gives an optimal interval of
about $\sqrt{2CM}$. \Cref{sec:pretrain:efficiency} returns to async checkpointing in the GPT-2 loop.
\end{systemslens}

% =====================================================================
\section{Experiment management}\label{sec:optim:explog}\covers{CS336 A1 \S5.3; CS336 A1 \S7.1}

You will run dozens of variants of every experiment, so the infrastructure for launching and comparing
them pays for itself within a day. \ident{llmsys/explog.py} implements the minimum we have found
useful:

\begin{itemize}
  \item \textbf{Typed configs with command-line overrides.} A \code{@dataclass} holds every
        hyperparameter; \code{config\_from\_cli(TrainConfig, sys.argv[1:])} applies
        \code{--key=value} overrides (and an optional \code{--config=file.json}) and \emph{rejects
        unknown keys}, so a typo in a sweep script fails immediately instead of silently running the
        default.
  \item \textbf{One directory per run, named by content.} \code{RunLogger.create(root, name, cfg)}
        appends a short hash of the config to the name and writes \code{config.json} and
        \code{meta.json} (command line, git revision, PyTorch version, GPU). Two runs with the same
        name but different settings cannot overwrite each other.
  \item \textbf{Append-only JSON-lines metrics, with step \emph{and} wall-clock time on every
        record.} JSONL survives crashes (at worst the last line is truncated, and the reader skips
        it), can be tailed live, and replotted against steps, tokens or seconds later. Plotting
        against wall-clock time is what tells you whether a change that improves loss per step is
        worth its cost per step.
  \item \textbf{Optional mirroring to Weights \& Biases.} If you pass \code{wandb\_project}, the
        same records are sent to W\&B; the local files remain the source of truth.
  \item \textbf{Sweeps as data.} \code{grid(lr=log\_spaced(1e-4, 1e-2, 5), batch=[32, 64])}
        produces points, and \code{sweep\_commands(...)} turns them into shell lines you can pipe
        into a Slurm array job or a cloud launcher.
\end{itemize}

\begin{lstlisting}[caption={A training script's experiment plumbing.},label={lst:optim:explog}]
@dataclass
class TrainConfig:
    d_model: int = 512; n_layers: int = 4; lr: float = 3e-3; batch: int = 64
    steps: int = 5000; warmup: int = 200; run_name: str = "base"

cfg = config_from_cli(TrainConfig, sys.argv[1:])          # --lr=1e-3 --batch=128
run = RunLogger.create("runs", cfg.run_name, cfg, wandb_project=None)
...
run.log(step, split="train", loss=loss.item(), lr=lr, grad_norm=gn,
        tokens=(step + 1) * cfg.batch * T, tok_per_s=tps)
if step % eval_every == 0:
    run.log(step, split="val", loss=val_loss)
run.summary(final_val=val_loss, diverged=False); run.close()
# later:  plot_runs(glob("runs/lr*"), y="loss", x="wall", where={"split": "val"})
\end{lstlisting}

\paragraph{What to log.} Every step: training loss, learning rate, pre-clip gradient norm, tokens
seen, throughput (tokens/s) and, on GPU, \code{torch.cuda.max\_memory\_allocated()}. Every few
hundred steps: validation loss on a \emph{fixed} set of batches (the same batches each time, so
curves are comparable), a few generated samples, and per-layer weight and update norms. Keep a
human-readable \emph{experiment log} next to the runs: what you changed, why, what you expected,
and what happened. Six weeks later the plots alone will not tell you why a run exists.

% =====================================================================
\section{Debugging the loop and tuning the learning rate}\label{sec:optim:debug}\covers{CS336 A1 \S7.2.3}

When a model does not learn, work from cheapest to most expensive check:

\begin{enumerate}
  \item \textbf{Initial loss $\approx \ln V$} (\cref{sec:optim:loss}).
  \item \textbf{Overfit one batch.} Train on a single fixed batch for a few hundred steps. The loss
        should approach zero. If it does not, the bug is in the model, loss or optimizer, not
        the data or hyperparameters. Typical culprits: a causal mask that leaks the
        future (loss goes to zero \emph{too} fast on the real data too), targets not shifted,
        parameters not registered (a plain Python list instead of \code{nn.ModuleList}),
        \code{optimizer.zero\_grad} never called.
  \item \textbf{Check gradients against finite differences} in float64 on a tiny model
        whenever you write a custom autograd function; PyTorch's \code{gradcheck} utility
        automates this.
  \item \textbf{Watch norms.} Log per-layer activation RMS, weight norms and the update-to-weight
        ratio $\|\Delta\theta\|/\|\theta\|$ (\code{llmsys.optim.update\_ratio}). For AdamW at a
        healthy learning rate it is around $10^{-3}$ per step for most matrices. Ratios near
        $10^{-1}$ mean the learning rate is far too high; ratios near $10^{-5}$ mean those weights
        barely move.
  \item \textbf{Sweep the learning rate on a log grid}, and judge the result by validation loss at a fixed
        token budget, with the schedule's $T_c$ set to that budget.
\end{enumerate}

\Cref{fig:optim:lrsweep} shows such a sweep, measured on CPU in this environment with Lab~9's
\code{run.py sweep-lr} on a synthetic first-order Markov language whose entropy rate (the lowest
achievable loss, 1.443 nats/token) is known exactly.

\begin{figure}[t]
\centering
\begin{tikzpicture}
\begin{groupplot}[group style={group size=2 by 1, horizontal sep=1.6cm}, width=0.5\linewidth, height=5.0cm,
  tick label style={font=\scriptsize}, label style={font=\small}, title style={font=\small}]
\nextgroupplot[title={(a) validation loss vs.\ step}, xlabel={step}, ylabel={val loss (nats/token)},
  ymin=1.4, ymax=3.6, xmin=0, xmax=400, legend style={font=\tiny, at={(0.98,0.98)}, anchor=north east}]
\addplot[black!40, dashed, forget plot] coordinates {(0,1.443) (400,1.443)};
\addplot[black!50, mark=none, line width=0.7pt] coordinates {(24,5.271) (49,4.167) (74,3.323) (99,2.777) (124,2.437) (149,2.218) (174,2.071) (199,1.974) (224,1.907) (249,1.859) (274,1.826) (299,1.801) (324,1.784) (349,1.771) (374,1.762) (399,1.754)};
\addlegendentry{$\eta=10^{-3}$}
\addplot[sysblue, line width=0.8pt] coordinates {(24,2.731) (49,1.898) (74,1.745) (99,1.744) (124,1.688) (149,1.659) (174,1.639) (199,1.618) (224,1.599) (249,1.569) (274,1.545) (299,1.528) (324,1.519) (349,1.499) (374,1.498) (399,1.49)};
\addlegendentry{$\eta=10^{-2}$}
\addplot[napkin, line width=0.8pt] coordinates {(24,2.264) (49,2.116) (74,2.213) (99,2.223) (124,2.151) (149,2.088) (174,2.096) (199,2.008) (224,2.0) (249,1.904) (274,1.847) (299,1.821) (324,1.737) (349,1.686) (374,1.665) (399,1.616)};
\addlegendentry{$\eta=0.3$}
\addplot[warnred, line width=0.8pt] coordinates {(24,3.437) (49,5.475) (74,4.936) (99,4.318) (124,4.175) (149,3.827) (174,3.296) (199,2.925) (224,2.838) (249,2.578) (274,2.467) (299,2.418) (324,2.286) (349,2.165) (374,2.074) (399,1.999)};
\addlegendentry{$\eta=1$}
\nextgroupplot[title={(b) final validation loss vs.\ $\eta$}, xlabel={peak learning rate $\eta$}, xmode=log,
  ymin=1.4, ymax=2.05, xmin=5e-4, xmax=2]
\addplot[black!40, dashed] coordinates {(5e-4,1.443) (2,1.443)};
\node[font=\tiny, black!60, anchor=south west] at (axis cs:6e-4,1.445) {entropy floor 1.443};
\addplot[sysblue, mark=*, line width=0.8pt] coordinates {(1e-3,1.754) (3e-3,1.512) (1e-2,1.490) (3e-2,1.490) (0.1,1.535) (0.3,1.616) (1,1.999)};
\end{groupplot}
\end{tikzpicture}
\caption{Learning-rate sweep for a 2-layer, $d=64$ Qwen3-shaped model trained for 400 steps
($16\times64$ tokens per step, warmup + cosine to $\eta/10$, AdamW, clip 1.0) on a synthetic
Markov language with a known entropy rate of 1.443 nats/token. Measured on CPU in this environment
(single seed, so differences under $\sim$0.01 are noise). The optimum is broad, spanning
$10^{-2}$ to $3\times10^{-2}$. At $\eta=1$ the training loss spiked above its initial value
(7.3 vs 5.8) and recovered only because clipping and QK-norm kept the run from diverging. The
best learning rate sits within roughly one decade of the largest one that still trains, as the
edge-of-stability folklore predicts.}
\label{fig:optim:lrsweep}
\end{figure}


\paragraph{Learning rate and batch size.} The two interact. Below a \emph{critical batch size}
(McCandlish et al., 2018), doubling the batch halves the gradient noise and lets you raise the
learning rate (roughly linearly for SGD, roughly as $\sqrt{B}$ for Adam) to reach the same loss in
half the steps. Above it, extra batch size buys almost nothing per token. The critical batch size grows as
the loss falls, which is one reason large runs ramp the batch size up during training. On one GPU,
the practical recipe is to use the largest batch that fits and still sits below the critical size,
re-tune the learning rate whenever you change it, and compare runs at equal \emph{tokens}, not
equal steps. \Cref{ch:scalinglaws} gives the scaling rules, including how $\mu$P (Yang et al.,
2022) makes the optimal learning rate transfer across model widths.

% =====================================================================
\section{Summary}\label{sec:optim:summary}
\begin{itemize}
  \item A training loop is a small system: sampler, numerically careful loss, optimizer with
        state, schedule, clipping, checkpointing and logging. It must be resumable and observable.
  \item Compute cross-entropy with logsumexp in fp32 and check that the initial loss is $\ln V$.
  \item Adam is RMSProp plus momentum plus bias correction. Its per-coordinate normalisation
        bounds the step by about $\eta$. AdamW decouples weight decay from that normalisation; for
        LLMs use $\beta_2\approx0.95$.
  \item The optimizer step is memory-bound: about 28 bytes per parameter per step for AdamW, a few
        milliseconds per billion parameters on an H100, as long as it runs as one fused or multi-tensor
        pass rather than hundreds of small kernels.
  \item AdamW model state is 16 bytes per parameter. Activations are linear in batch size, and attention
        scores dominate them without FlashAttention. Derive $\text{mem}(B)=aB+b$ before every run.
  \item Warmup protects Adam's early steps; cosine bakes the run length into the schedule; WSD
        does not, which allows branching decays.
  \item Clip the \emph{global} gradient norm, log the pre-clip norm, and remember that clipping is a
        global reduction under sharding.
  \item Lion, Adafactor, Shampoo/SOAP and Muon trade memory or extra compute for better
        conditioning. Muon orthogonalises momentum with a few tensor-core matmuls.
  \item Sample random windows from a validated \code{uint16} memmap; checkpoint weights,
        optimizer state, step, sampler state and RNG atomically; test resumption bit for bit.
  \item Log step and wall-clock time on every record, keep configs typed and runs uniquely named,
        and debug from the cheapest check (initial loss) to the most expensive (LR sweeps).
\end{itemize}

% =====================================================================
\section{Exercises}\label{sec:optim:exercises}

\begin{exercise}[Bias correction]
Show that without bias correction, Adam's first update with $\beta_1=0.9,\beta_2=0.999$ has
magnitude $\eta(1-\beta_1)/\sqrt{1-\beta_2}\approx3.16\,\eta$ per coordinate. What happens with
$\beta_2=0.95$? Why is the second-moment correction the more important one?
\end{exercise}

\begin{exercise}[Coupled versus decoupled decay]
Consider a single scalar parameter with zero loss gradient, trained with Adam + L2 ($\lambda$)
and with AdamW ($\lambda$), both at learning rate $\eta$. Write the update for each. Show that
Adam + L2 shrinks the weight by $\approx\eta\,\mathrm{sign}(\theta)$ per step (independent of
$\lambda$, once $v$ is dominated by the decay term) while AdamW shrinks it geometrically.
\end{exercise}

\begin{exercise}[Optimizer bytes]
Count bytes read and written per parameter for one step of (a) SGD with momentum in fp32,
(b) Lion in mixed precision with fp32 master weights, (c) AdamW with bf16 $m$ and $v$ and fp32
master weights. Convert each into milliseconds per billion parameters on an H100 (3.35\,TB/s).
\end{exercise}

\begin{exercise}[Launch-bound updates]
A 28-layer Qwen3-0.6B has about 310 parameter tensors. If each elementwise kernel launch costs
5\,$\mu$s of CPU time and the loop implementation issues 10 kernels per tensor, how long does
the optimizer step take on the CPU side? Compare with the memory-bound time for 0.6B parameters.
Measure the loop and \code{foreach} paths of \code{llmsys.optim.AdamW} on a GPU if you have one.
\end{exercise}

\begin{exercise}[Accounting method]
Using the method of \cref{sec:optim:account}, derive $\text{mem}(B)=aB+b$ for Qwen3-0.6B
(\cref{ch:modern}) trained in mixed precision with AdamW at $T=4096$, (i) with naive attention and
(ii) with FlashAttention. What micro-batch size fits on a 24\,GB GPU in each case?
\end{exercise}

\begin{exercise}[Schedules and resumption]
A colleague resumes a cosine run from a checkpoint, but their script reconstructs the schedule
with \code{total} equal to the \emph{remaining} steps and starts counting at zero. Sketch the
learning rate they actually used. What would the loss curve look like?
\end{exercise}

\begin{exercise}[Clipping semantics under sharding]
Two tensor-parallel ranks each hold half of every gradient. Rank 0 has local norm 3 and rank 1
local norm 4; $M=1$. What scale should each rank apply? What happens if each rank clips its
own shard to norm 1 independently? Write the extra collective needed.
\end{exercise}

\begin{exercise}[Muon geometry]
Show that $p(X)=aX+b(XX^\top)X+c(XX^\top)^2X$ maps $U\Sigma V^\top$ to $Up(\Sigma)V^\top$. Why must
the input be scaled so that $\sigma_{\max}\le1$? Why does Muon apply only to 2-D hidden matrices
and not to embeddings?
\end{exercise}

\begin{exercise}[Resumable data]
Our sampler draws with replacement. Design a sampler that visits every window of a token file
exactly once per epoch in random order without materialising the permutation (hint: a random
affine bijection $i\mapsto (ai+c)\bmod n$ with $\gcd(a,n)=1$). What must its \code{state\_dict}
contain?
\end{exercise}

\begin{exercise}[Edge of stability]
Reproduce \cref{fig:optim:lrsweep} with \code{run.py sweep-lr}, extend the grid until a run
diverges without clipping (pass \code{--clip=1e9}), and plot the final loss against the
learning rate. How far below the divergence point is the best learning rate? Repeat after
removing QK-norm from the model. What changes?
\end{exercise}

% =====================================================================
\section{Lab 9: Optimizers and a resumable training loop}\label{sec:optim:lab}

\begin{labbox}{9}{Optimizers and a resumable training loop}
\labdir{lab09_optim}
\labsection{Goals}
Implement the update rules and the plumbing of this chapter, verify them against PyTorch, and
prove that your loop can be killed and resumed without changing a single bit of the result.

\labsection{Setup}
\begin{lstlisting}[style=shell]
cd code
LAB_IMPL=starter  python -m pytest labs/lab09_optim -q     # your implementation
LAB_IMPL=solution python -m pytest labs/lab09_optim -q     # reference
python labs/lab09_optim/run.py --help                       # experiments
\end{lstlisting}

\labsection{Tasks}
\begin{enumerate}
  \item \code{SGD.step}: momentum and Nesterov, matching \code{torch.optim.SGD} (4 tests).
  \item \code{AdamW.\_step\_loop}: decoupled and coupled decay, bias correction, $\epsilon$
        placement; matches \code{torch.optim.AdamW} and \code{torch.optim.Adam}.
  \item \code{AdamW.\_step\_foreach}: the same numbers with \code{torch.\_foreach\_*} kernels.
  \item \code{WarmupCosine} and \code{WSD} schedules (0-based steps, floor after the end).
  \item \code{clip\_grad\_norm\_}: global norm in fp32, in-place scaling, pre-clip norm returned;
        matches \code{torch.nn.utils.clip\_grad\_norm\_}.
  \item \code{RandomWindowBatcher}: vectorised window gather from a memmap, shifted targets,
        \code{state\_dict}/\code{load\_state\_dict} for exact replay.
  \item \code{newton\_schulz\_orthogonalize}: singular values pulled into $[0.5,1.3]$.
  \item \code{optimizer\_state\_bytes} and \code{training\_state\_bytes}.
  \item \code{RunLogger.log} and \code{load\_jsonl} (tolerating a truncated last line).
  \item \code{CheckpointManager.save/restore}: atomic write, RNG and data state, keep-last-$k$.
  \item \code{train\_loop}: wire everything together; the bit-exact resume test must pass.
\end{enumerate}

\labsection{Measurements}
\begin{itemize}
  \item \code{run.py make-markov} then \code{run.py sweep-lr}: reproduce
        \cref{fig:optim:lrsweep}; extend the grid to find the divergence point with and without
        clipping.
  \item \code{run.py compare}: tune each optimizer's learning rate with a three-point sweep
        before comparing. Does the ranking survive tuning?
  \item \code{run.py wsd}: compare cosine runs and WSD branches at three budgets
        (\cref{fig:optim:wsd}). Count total optimizer steps for both approaches.
  \item On a GPU: time \code{AdamW(foreach=False)}, \code{AdamW(foreach=True)} and
        \code{torch.optim.AdamW(fused=True)} for Qwen3-0.6B's parameters; compare with the
        28-bytes-per-parameter bound.
\end{itemize}

\labsection{Deliverables}
Passing tests; the three plots above with one paragraph each; a table of measured optimizer-step
times against the bandwidth bound.

\labsection{Stretch goals}
Implement 8-bit block-wise quantised Adam states and measure the loss difference after 1000
steps. Add a \code{PrefetchBatcher} to \code{train\_loop} and measure the step-time change on a
GPU. Implement Adafactor's factored second moment and compare its memory.
\end{labbox}
